
\documentclass[11pt,a4paper]{article}
\usepackage[T1]{fontenc}
\usepackage{lmodern}
\usepackage{amsmath,amssymb,amsfonts,mathtools}
\usepackage{graphicx,booktabs,longtable,array,tabularx,adjustbox}
\usepackage{geometry,microtype,fancyhdr,lastpage,etoolbox,enumitem}
\usepackage{newunicodechar,textcomp}
\usepackage{hyperref}
\hypersetup{hidelinks,pdfcreator={XeLaTeX / QM1 academic edition}}
\geometry{left=20mm,right=20mm,top=22mm,bottom=23mm}
\setcounter{secnumdepth}{2}
\setlength{\parindent}{0pt}
\setlength{\parskip}{4pt plus 1pt}
\setlength{\emergencystretch}{2em}
\setlength{\tabcolsep}{3pt}
\renewcommand{\arraystretch}{1.16}
\setlist{nosep,topsep=3pt,itemsep=2pt,leftmargin=7mm}
\AtBeginEnvironment{longtable}{\footnotesize\setlength{\tabcolsep}{2.2pt}}
\newcommand{\tightlist}{\setlength{\itemsep}{0pt}\setlength{\parskip}{0pt}}
\providecommand{\Cov}{\operatorname{Cov}}
\providecommand{\Var}{\operatorname{Var}}
\providecommand{\Corr}{\operatorname{Corr}}
\providecommand{\MAD}{\operatorname{MAD}}
\providecommand{\E}{\mathbb{E}}
\providecommand{\R}{\mathbb{R}}
\newunicodechar{—}{\textemdash}
\newunicodechar{–}{\textendash}
\newunicodechar{€}{\texteuro}
\newunicodechar{≤}{\ensuremath{\leq}}
\newunicodechar{≥}{\ensuremath{\geq}}
\newunicodechar{×}{\ensuremath{\times}}
\newunicodechar{÷}{\ensuremath{\div}}
\newunicodechar{Σ}{\ensuremath{\Sigma}}
\newunicodechar{σ}{\ensuremath{\sigma}}
\newunicodechar{μ}{\ensuremath{\mu}}
\newunicodechar{Δ}{\ensuremath{\Delta}}
\newunicodechar{∑}{\ensuremath{\sum}}
\newunicodechar{≈}{\ensuremath{\approx}}
\newunicodechar{√}{\ensuremath{\surd}}
\newunicodechar{≠}{\ensuremath{\ne}}
\newunicodechar{→}{\ensuremath{\to}}
\newunicodechar{∞}{\ensuremath{\infty}}
\newunicodechar{∈}{\ensuremath{\in}}
\newunicodechar{−}{\ensuremath{-}}
\newunicodechar{²}{\ensuremath{^{2}}}
\newunicodechar{³}{\ensuremath{^{3}}}
\newunicodechar{₀}{\ensuremath{_{0}}}
\newunicodechar{₁}{\ensuremath{_{1}}}
\newunicodechar{₂}{\ensuremath{_{2}}}
\newunicodechar{₃}{\ensuremath{_{3}}}
\newunicodechar{₄}{\ensuremath{_{4}}}
\newunicodechar{ₜ}{\ensuremath{_{t}}}
\newunicodechar{ₚ}{\ensuremath{_{p}}}
\newunicodechar{ₙ}{\ensuremath{_{n}}}
\newunicodechar{ₓ}{\ensuremath{_{x}}}
\newunicodechar{ᵢ}{\ensuremath{_{i}}}
\newunicodechar{ŷ}{\ensuremath{\hat y}}
\newunicodechar{Ŷ}{\ensuremath{\widehat Y}}
\newunicodechar{ₘ}{\ensuremath{_{m}}}
\newunicodechar{ₑ}{\ensuremath{_{e}}}
\newunicodechar{ₐ}{\ensuremath{_{a}}}
\newunicodechar{ₖ}{\ensuremath{_{k}}}
\newunicodechar{ₛ}{\ensuremath{_{s}}}
\newunicodechar{ₗ}{\ensuremath{_{l}}}
\newunicodechar{ₕ}{\ensuremath{_{h}}}
\newunicodechar{ₒ}{\ensuremath{_{o}}}
\newunicodechar{₊}{\ensuremath{_{+}}}
\newunicodechar{₋}{\ensuremath{_{-}}}
\newunicodechar{₌}{\ensuremath{_{=}}}
\newunicodechar{₍}{\ensuremath{_{(}}}
\newunicodechar{₎}{\ensuremath{_{)}}}
\newunicodechar{⁻}{\ensuremath{^{-}}}
\newunicodechar{⁺}{\ensuremath{^{+}}}
\newunicodechar{⁰}{\ensuremath{^{0}}}
\newunicodechar{¹}{\ensuremath{^{1}}}
\newunicodechar{⁴}{\ensuremath{^{4}}}
\newunicodechar{⁵}{\ensuremath{^{5}}}
\newunicodechar{⁶}{\ensuremath{^{6}}}
\newunicodechar{⁷}{\ensuremath{^{7}}}
\newunicodechar{⁸}{\ensuremath{^{8}}}
\newunicodechar{⁹}{\ensuremath{^{9}}}
\makeatletter
\renewcommand{\section}{\@startsection{section}{1}{0pt}{2.0ex plus .7ex minus .3ex}{1.0ex}{\large\bfseries}}
\renewcommand{\subsection}{\@startsection{subsection}{2}{0pt}{1.8ex plus .5ex minus .2ex}{.7ex}{\normalsize\bfseries}}
\renewcommand{\subsubsection}{\@startsection{subsubsection}{3}{0pt}{1.3ex plus .4ex minus .2ex}{.5ex}{\normalsize\bfseries}}
\makeatother
\pagestyle{fancy}
\fancyhf{}
\renewcommand{\headrulewidth}{0.35pt}
\renewcommand{\footrulewidth}{0pt}
\fancyhead[L]{\small\textsc{KEDGE Business School --- QM1}}
\fancyhead[R]{\small\textsc{2026--2027}}
\fancyfoot[L]{\small Badr ABOUFARASSE, PhD}
\fancyfoot[R]{\small\thepage\ /\ \pageref{LastPage}}
\setlength{\headheight}{14pt}
\setkeys{Gin}{keepaspectratio}
\begin{document}

\begin{center}
{\small\textsc{KEDGE Business School \quad EBP1 \quad Quantitative Methods 1}}\\[6pt]
{\Large\bfseries Fiche 4 — Distributions bivariées et régression}\\[6pt]
{\normalsize\bfseries Énoncés des exercices}\\[6pt]
{\small Badr ABOUFARASSE, PhD \quad --- \quad 2026--2027}
\end{center}
\vspace{2mm}\hrule\vspace{6mm}

\hypertarget{applications-fondamentales}{%
\section{Applications fondamentales}\label{applications-fondamentales}}

Maîtriser les définitions et les calculs essentiels du cours.

\emph{Application du cours --- 15 min}

\hypertarget{exercice-3-droite-de-mayer}{%
\subsection{Exercice 3 -- Droite de Mayer}\label{exercice-3-droite-de-mayer}}

Pour les points {\((1,7),(2,10),(3,10),(4,14),(5,14)\)} : former deux sous-nuages, calculer leurs points moyens {\(A\)} et {\(B\)}, puis déterminer la droite de Mayer.

\hypertarget{pruxe9paration-aux-uxe9valuations}{%
\section{Préparation aux évaluations}\label{pruxe9paration-aux-uxe9valuations}}

Résoudre des exercices de contrôle continu et de partiel en autonomie.

\emph{Préparation CC / partiel --- 35 min}

\hypertarget{exercice-1-ruxe9gression-linuxe9aire-compluxe8te}{%
\subsection{Exercice 1 -- Régression linéaire complète}\label{exercice-1-ruxe9gression-linuxe9aire-compluxe8te}}

On observe : {\[(1,3),\ (2,5),\ (3,5),\
(4,8),\ (5,9).\]}

\begin{enumerate}
\item
  Tracer le nuage de points et décrire direction, forme et intensité apparente.
\item
  Calculer {\(\bar x\)}, {\(\bar y\)}, {\(V(x)\)}, {\(V(y)\)}, {\(\Cov(x,y)\)} et le coefficient de corrélation {\(r\)}.
\item
  Estimer la droite des moindres carrés {\(\hat y=ax+b\)}.
\item
  Calculer les valeurs ajustées et les résidus {\(e_i=y_i-\hat y_i\)}.
\item
  Calculer {\(SSE\)}, {\(SST\)} et {\(R^2\)} ; interpréter {\(R^2\)}.
\end{enumerate}

\emph{Préparation CC / partiel --- 15 min}

\hypertarget{exercice-4-pruxe9vision-en-deux-uxe9tapes}{%
\subsection{Exercice 4 -- Prévision en deux étapes}\label{exercice-4-pruxe9vision-en-deux-uxe9tapes}}

\begin{longtable}[]{@{}lllllllll@{}}
\toprule\noalign{}
Mois & 1 & 2 & 3 & 4 & 5 & 6 & 7 & 8 \\
\midrule\noalign{}
\endhead
\bottomrule\noalign{}
\endlastfoot
Publicité {\(X\)} & 12 & 15 & 14 & 18 & 20 & 22 & 25 & 27 \\
Ventes {\(Y\)} & 48 & 52 & 50 & 59 & 63 & 68 & 73 & 78 \\
\end{longtable}

On donne les deux ajustements : {\[\widehat{Y}=2.064X+21.899,
\qquad
\widehat{X}_t=2.155t+9.429.\]} Prévoir la dépense publicitaire puis les ventes au mois 10. Identifier le caractère d'extrapolation de la première étape.

\emph{Préparation CC / partiel --- 45 min}

\hypertarget{exercice-5-deux-ruxe9gressions-successives-type-partiel}{%
\subsection{Exercice 5 --- Deux régressions successives --- type partiel}\label{exercice-5-deux-ruxe9gressions-successives-type-partiel}}

Sur six périodes, on observe le rang t, la publicité x et les ventes y.

\begin{longtable}[]{@{}lllllll@{}}
\toprule\noalign{}
t & 1 & 2 & 3 & 4 & 5 & 6 \\
\midrule\noalign{}
\endhead
\bottomrule\noalign{}
\endlastfoot
x & 4 & 5 & 7 & 8 & 10 & 11 \\
y & 10 & 12 & 15 & 18 & 22 & 24 \\
\end{longtable}

\begin{enumerate}
\tightlist
\item
  Construire le tableau complet t, x, y, t², x², y², tx, xy et les sommes.
\item
  Calculer les moyennes, variances, covariance, r et la droite de x en fonction de t.
\item
  Faire de même pour y en fonction de x.
\item
  Prévoir x puis y pour t=8 ; qualifier les extrapolations.
\end{enumerate}

\hypertarget{approfondissement-facultatif}{%
\section{Approfondissement facultatif}\label{approfondissement-facultatif}}

Aller plus loin dans l'interprétation et les applications professionnelles.

\emph{Approfondissement --- 20 min}

\hypertarget{exercice-2-interpolation-extrapolation-causalituxe9-et-influence}{%
\subsection{Exercice 2 -- Interpolation, extrapolation, causalité et influence}\label{exercice-2-interpolation-extrapolation-causalituxe9-et-influence}}

À partir du modèle de l'exercice 1 :

\begin{enumerate}
\item
  prévoir {\(Y\)} pour {\(X=3.5\)} et qualifier la prévision ;
\item
  prévoir {\(Y\)} pour {\(X=10\)} et expliquer pourquoi cette prévision est risquée ;
\item
  ajouter le point {\((10,30)\)} et recalculer la pente et l'ordonnée à l'origine ; commenter l'influence de ce point ;
\item
  expliquer pourquoi une forte corrélation ne prouve pas une relation causale.
\end{enumerate}

\emph{Approfondissement --- 20 min}

\hypertarget{compluxe9ment-1-diagnostic-de-ruxe9gression-et-duxe9cision-responsable}{%
\subsection{Complément 1 --- Diagnostic de régression et décision responsable}\label{compluxe9ment-1-diagnostic-de-ruxe9gression-et-duxe9cision-responsable}}

Un modèle reliant fréquentation et température donne r = 0,78 et R² = 0,61. Le graphique des résidus est en U, un week-end d'été présente un résidu très élevé et le responsable veut prévoir une température supérieure de 12 °C au maximum observé.

\begin{enumerate}
\tightlist
\item
  Interprétez r et R².
\item
  Diagnostiquez les résidus.
\item
  Discutez le week-end atypique et la prévision.
\item
  Recommandez les étapes suivantes.
\end{enumerate}

\textbf{Indication méthodologique :}

Distinguez intensité, forme fonctionnelle, observation inhabituelle, extrapolation et causalité.

\end{document}
